\scanpage{25}\mspage{III}
\margin{\tikz[x=.08cm,y=.08cm]{\draw (0,0)..controls (1,1.1) and (2.1,1.9)..(1.4,2.8);}}%
\uline{Lemme I}. Soit $\Gamma$ un $J$-polygone et $x,y\in\Gamma$ \inserted{$x\ne y$} ; si $\Gamma^{*}_{xy}$ est l'un de deux arcs de polygone réliant $x$ à $y$, décrit dans la Remarque 1, \margin{!\ \tikz[x=.07cm,y=.07cm]{\draw (0,0)..controls (.4,2) and (1.7,2.4)..(1.5,3.4);}}il \uline{existe un système} fondamental \inserted{$\beta$} de voisinages \deleted{$\unsure{\beta|B}$} \inserted{ouverts} de $\Gamma^{*}_{xy}$ tel que $V\in\beta$ vérifie :
\margin{$\cdot$}
\[
\begin{aligned}
&V\cap\mathrm{int}\,\Gamma\quad\text{et}\quad V\cap\mathrm{ext}\,\Gamma\quad\text{sont connexes (donc}\\
&V-V\cap\Gamma=(V\cap\mathrm{int}\,\Gamma)\cup(V\cap\mathrm{ext}\,\Gamma)\ \text{a deux composantes connexes)}
\end{aligned}
\]

\margin{!}\uline{Demonstration}. gr\uline{âce à l'axiom du c}hoix il suffit de prouver que si $U$ est un voisinage \uline{de $\Gamma$}, \margin{!\\$\Gamma_{xy}$}il existe un voisinage ouvert $V\subset U$ satisfaisant \inserted{à} la condition de l'ennoncé. Pour chaque $z\in\Gamma_{xy}$ on choisit un disque \inserted{ouvert} $D_{z}\subset U$ et satisfaisant b) de la definition 1\unsure{3|,|;} par rapport a $\Gamma$. Il est clair que $\{D_{z}\}_{\inserted{z\in\Gamma_{xy}}}$ est un récouvrement ouvert de l'ensemble compact $\Gamma_{xy}$. Soit donc un sous-récouvrement fini $\{D_{z_{i}}\}^{n}_{i=1}$ de $\Gamma_{xy}$.

Soit $I_{i}$ l'intervalle \inserted{ouvert ou semiouvert de $\Gamma_{xy}$} decoupé par $D_{z_{i}}$, i.e. \inserted{\tikz[baseline=(f.base)]{\node[inner xsep=.12em,inner ysep=.12ex] (f) {$\deleted{\unsure{I|F}}\Gamma_{xy}\cap D_{z_{i}}=I_{i}.$}; \draw[line width=.5pt] ([xshift=-.65em]f.south west)--(f.south west)--(f.north west); \draw[line width=.5pt] ([xshift=.30em,yshift=.35ex]f.north west)--([yshift=.65ex]f.north east)..controls ([xshift=.12em,yshift=.58ex]f.north east) and ([xshift=.15em]f.south east)..([xshift=.23em]f.south east)--([xshift=.45em]f.south east);}} On observe que si $I_{i}\cap I_{j}\ne\phi$ alors $D_{z_{i}}\cap D_{z_{j}}\cap\mathrm{int}\,\Gamma\ne\phi$. Cela permet de montrer que si $z',z''\in\bigcup\limits_{i=1}^{n}D_{z_{i}}\cap\mathrm{int}\,\Gamma$ il existe une suite :

$D_{z_{i_{1}}}\cap\mathrm{int}_{\Gamma},\ \ldotp\;\ldotp\;-\,{-},\ D_{z_{i_{m}}}\cap\mathrm{int}\,\Gamma$ d'ouverts connexes (par le choix de chaque $D_{z}$) tels que $D_{z_{i_{\alpha}}}\cap D_{z_{i_{\alpha+1}}}\cap\mathrm{int}\,\Gamma\ne\phi$ pour $1\le\alpha\le m-1$, et $z'\in D_{z_{i_{1}}}\cap\mathrm{int}_{\Gamma},\ z''\in D_{z_{i_{m}}}\cap\mathrm{int}\,\Gamma$. Ces deux propriétés verifiées par $\bigcup\limits_{i=1}^{n}D_{z_{i}}\cap\mathrm{int}\,\Gamma$ suffisent à prouver la connexité de $\bigcup\limits_{i=1}^{n}D_{z_{i}}\cap\mathrm{int}\,\Gamma$.

\margin{!\ $\cdot$}\deleted{\unsure{Une|Un}} Le même argument prouve (qui\uline{tte à changer $\mathrm{int}\,\Gamma$ par $\mathrm{ext}\,\Gamma$}) que $\bigcup\limits_{i=1}^{n}D_{z_{i}}\cap\mathrm{ext}\,\Gamma$ est connexe. Par conséquent il suffit de prendre $V=\bigcup\limits_{i=1}^{n}D_{z_{i}}$, Cela montre le lemme.

\margin{$\bullet$\quad\uline{?}}On\textsuperscript{\unsure{$*$}} utilise le fait, \uline{sans le dire}, que \deleted{tout couple de} \inserted{deux} points $x,y$ \deleted{\unsure{de|un}} \inserted{d'un} ouvert connexe de $\mathbb{R}^{2}$ peu\deleted{t}\inserted{vent} être reliés par un arc de polygone.
% archive sheet number: 24
